\documentclass[11pt]{article}
\usepackage[margin=1.1in]{geometry}
\usepackage{amsmath,amssymb}
\usepackage{booktabs}
\usepackage{tikz}
\usetikzlibrary{positioning,fit,backgrounds}
\usepackage[hidelinks]{hyperref}
\usepackage{parskip}

\newcommand{\CNOT}{\textsc{cnot}}
\newcommand{\CCNOT}{\textsc{ccnot}}
\newcommand{\code}[1]{\texttt{#1}}

\title{\textbf{The Sliced Sandwich Construction}\\[2pt]
\large A two-sided keyed slicing mechanism on $2n$ wires}
\author{local\_mixing project}
\date{July 20, 2026; balanced variant September 4, 2026 \\
\small branch \code{ssg-gen-mix-clean}}

\begin{document}
\maketitle

\begin{abstract}
\noindent The \emph{sliced sandwich} is a keyed obfuscation construction that
computes a source function $C$ on one designated input slice and, on the same
slice, reveals only a fresh random function $D^{-1}$ when run backwards. It is
a slicing mechanism specific to the sandwich structure and is distinct from
the gadget and feistel paths: it uses no secret sharing, placing $C$ and a
random padding computation $D$ as plaintext g57 circuits on the low half of
$2n$ wires, with two slice blocks of \CNOT/\CCNOT{} gates interleaved through
the two computations. It comes in two variants, selected by a choice bit: the
original \emph{classic} layout, in which both computations share the low half
and both directions read out on the high one, and a \emph{balanced} layout, in
which each half hosts one computation and each direction has its own slice and
its own answer half. This note gives both constructions, their slice semantics
in both directions, the rationale, and the parameters.
\end{abstract}

\section{Wires and structure}

The transformed circuit acts on $2n$ wires in two halves of $n$:
the first half (wires $0..n{-}1$) holds the input $x$, and the second half
(wires $n..2n{-}1$) holds the auxiliary $y$ --- the \emph{slice register}.

The circuit is a sandwich of two computations around a copy step:
\[
A \;=\; \bigl[\, C \parallel S_1 \,\bigr]\;;\;\; N \;;\;\; \bigl[\, D \parallel S_2 \,\bigr],
\]
where $\parallel$ denotes a random interleaving. The pieces:

\begin{itemize}
\item \textbf{$C$} --- the source circuit, a g57 circuit placed directly on
wires $0..n{-}1$. No secret sharing: this is a plaintext computation whose
obfuscation comes from the subsequent mixing.
\item \textbf{$D$} --- a fresh random circuit of $m$ g57 gates on wires
$0..n{-}1$, the \emph{same design} as the $C$ block (target and two distinct
controls per gate). Default $m = n(\log_2 n)^2$.
\item \textbf{$N$} --- the copy step $y \mathrel{\hat{}}= x$: the $n$ \CNOT{}s
$\;\text{wire } (n{+}i) \mathrel{\hat{}}= \text{wire } i$. (This is the
classic direction; see \S\ref{sec:balanced}.)
\item \textbf{$S_1, S_2$} --- two independent slice blocks of $s$ gates each
(default $s = n\log_2 n$). Every gate targets the first half and reads at
least one second-half wire with positive polarity:
\[
\underbrace{x_i \mathrel{\hat{}}= y_j}_{\text{\CNOT, }\approx 1/3}
\qquad
\underbrace{x_i \mathrel{\hat{}}= x_j \wedge y_k}_{\text{\CCNOT, rest}}
\]
so each block is \emph{dead when the second half is zero}. $S_1$ is randomly
interleaved with $C$, $S_2$ with $D$; the interleaving is uniform and
preserves each computation's internal order (necessary for $C$ and $D$ to
compute correctly).
\end{itemize}

\section{Slice semantics (classic)}

\subsection{Forward: the slice reveals $C$}

On the zero slice $y=0$, trace the running state $(a,b)$ from $(x,0)$:
\[
(x,0)
\xrightarrow{\;C \parallel S_1\;} (C(x),\,0)
\xrightarrow{\;N\;} (C(x),\,C(x))
\xrightarrow{\;D \parallel S_2\;} (\text{junk},\,C(x)).
\]
During the first stage the second half is still $0$, so every $S_1$ gate is
dead and $C$ computes correctly. The copy step moves $C(x)$ into the second
half. In the third stage $S_2$ is \emph{live} (the second half now holds
$C(x)\neq 0$), but it only targets the first --- junk --- half, so it cannot
corrupt the answer. Hence
\[
\boxed{\,A(x,0) = (\text{junk},\, C(x))\,}
\]
with the answer on the second half.

\subsection{Inverse: the same slice reveals $D^{-1}$}

Every gate is an involution, so $A^{-1}$ is the reversed gate list. On input
$(p,0)$:
\[
(p,0)
\xrightarrow{\;(D \parallel S_2)^{-1}\;} (D^{-1}(p),\,0)
\xrightarrow{\;N\;} (D^{-1}(p),\,D^{-1}(p))
\xrightarrow{\;(C \parallel S_1)^{-1}\;} (\text{junk},\,D^{-1}(p)).
\]
Now it is $S_2$ that is dead (its stage runs first, second half $0$), so
$D^{-1}$ computes cleanly; and $S_1$ that fires (its stage sees a nonzero
second half), junking the first half. Hence
\[
\boxed{\,A^{-1}(p,0) = (\text{junk},\, D^{-1}(p)).\,}
\]

The construction is thus sliced on the \emph{same} slice in both directions:
each slice block guards one direction's correctness ($S_1$ forward, $S_2$
inverse) and scrambles the other direction's junk. The forward oracle gives
$C$; the backward oracle gives only the random $D^{-1}$, never $C^{-1}$.

\subsection{Off-slice}

For $y \neq 0$ the second-half output is $y \oplus C'(x,y)$, where $C'$ is the
$S_1$-disturbed computation. The $y\oplus$ masking from the copy step
guarantees the output differs from the clean answer $C(x)$ on every nonzero
slice, independently of whether the disturbance $C'$ happens to coincide with
$C$ at a given $x$.

\section{The balanced variant}
\label{sec:balanced}

The \emph{balanced} variant changes two things and leaves everything else ---
$C$, $S_1$, the interleaving, the float stage --- untouched:

\begin{itemize}
\item the copy step flips direction, $N: x \mathrel{\hat{}}= y$
      (\;$\text{wire } i \mathrel{\hat{}}= \text{wire } (n{+}i)$\;), and
\item the $D$ block moves to the second half, $D'$ on wires $n..2n{-}1$.
\end{itemize}

\noindent A third change is forced by those two: $S_2$ must \emph{mirror} with
$D$, targeting the second half and reading the first,
\[
\underbrace{y_i \mathrel{\hat{}}= x_j}_{\text{\CNOT, }\approx 1/3}
\qquad
\underbrace{y_i \mathrel{\hat{}}= y_j \wedge x_k}_{\text{\CCNOT, rest}}
\]
so that $S_2'$ is dead when the \emph{first} half is zero. Otherwise $S_2$
would fire as soon as $D'$ makes the second half nonzero, and would overwrite
the forward answer, which in this layout stays in the first half. So
\[
A \;=\; \bigl[\, C \parallel S_1 \,\bigr]\;;\;\; N(x \mathrel{\hat{}}= y) \;;\;\;
        \bigl[\, D' \parallel S_2' \,\bigr].
\]

\subsection{Forward: the slice reveals $C$ on the first half}

On the same zero slice $y=0$, from $(x,0)$:
\[
(x,0)
\xrightarrow{\;C \parallel S_1\;} (C(x),\,0)
\xrightarrow{\;N\;} (C(x),\,0)
\xrightarrow{\;D' \parallel S_2'\;} (C(x),\,\text{junk}).
\]
Block 1 is unchanged, so $S_1$ is dead and $C$ computes correctly; the flipped
copy step is then a no-op on the slice; and block 2 writes \emph{only} the
second half, so the answer is frozen where block 1 left it:
\[
\boxed{\,A(x,0) = (C(x),\,\text{junk})\,}
\]
with the answer on the first half. Where the classic variant protects the
answer by moving it into a register that the second block never targets, the
balanced variant protects it by leaving it in a half that the second block
never targets.

\subsection{Inverse: the mirrored slice reveals $D^{-1}$}

The inverse is sliced at $x=0$, not $y=0$. On input $(0,q)$:
\[
(0,q)
\xrightarrow{\;(D' \parallel S_2')^{-1}\;} (0,\,D^{-1}(q))
\xrightarrow{\;N\;} (D^{-1}(q),\,D^{-1}(q))
\xrightarrow{\;(C \parallel S_1)^{-1}\;} (\text{junk},\,D^{-1}(q)).
\]
$S_2'$ is dead because the first half is zero, so $D^{-1}$ computes cleanly on
the second half; $N$ copies it down; and the reversed block 1 --- $C^{-1}$ with
$S_1$ now firing on a nonzero second half --- junks the first half only. Hence
\[
\boxed{\,A^{-1}(0,q) = (\text{junk},\, D^{-1}(q)).\,}
\]

\subsection{Off-slice}

For $y \neq 0$ the first-half output is $C'(x,y) \oplus y$ with $C'$ the
$S_1$-disturbed computation: the flipped copy step supplies exactly the same
$y\oplus$ masking as the classic one, now on the answer half, so no nonzero
slice yields the clean answer.

\subsection{What is balanced, and what it costs}

Each half now hosts one computation and is one direction's payload: the
classic layout's low half is a shared workspace and its high half a pure
answer register, whereas here $C$ owns the first half and $D$ the second.
The two directions consequently have \emph{different} slices ($y=0$ forward,
$x=0$ backward), where the classic construction slices both at $y=0$.

Equivalently: the classic layout has a half --- the low one --- that is junk
in \emph{both} directions, and the balanced layout has none. One asymmetry
comes with it: on the forward slice the $D'$ lane starts from the constant $0$
rather than from live data, so its junk output is a function of $x$ only
through the $S_2'$ injections, whereas the classic layout's junk half starts
from $C(x)$ and runs $D$ on it.

\section{Rationale}

\textbf{Why interleave rather than prepend.} In the plain gadget path the
slice block sits at the input port. Here the sandwich has two computations and
a natural two-sided structure, and weaving each slice block \emph{through} its
computation ties the disturbance to intermediate states rather than only the
boundary, which the mixing then disperses. Correctness is unaffected because
the slice gates are dead on the slice regardless of their positions.

\textbf{Why $D$ mirrors $C$.} $D$ is drawn from the same distribution as the
source blocks (random g57 gates on the low $n$ wires) so that the forward and
inverse directions are structurally indistinguishable: an adversary holding
the mixed circuit sees two same-design computations and cannot tell from gate
statistics which end carries the ``real'' function.

\textbf{Why this gate vocabulary.} The slice gates are positive-polarity
\CNOT/\CCNOT{} --- the same mixing-native vocabulary as the rest of the
circuit, with no complemented gates and no polarity pattern that could betray
the slice. The slice is the canonical zero vector, so nothing about it is
encoded in the circuit beyond the deadness structure.

\section{Parameters}

\begin{center}
\renewcommand{\arraystretch}{1.2}
\begin{tabular}{clc}
\toprule
symbol & meaning & default \\
\midrule
$n$ & wires of the source function (circuit is $2n$ wires) & --- \\
$m$ & gates in the random $D$ computation & $n(\log_2 n)^2$ \\
$s$ & gates in each slice block $S_1, S_2$ & $n\log_2 n$ \\
variant & \code{classic} or \code{balanced} (\S\ref{sec:balanced}) & \code{classic} \\
\bottomrule
\end{tabular}
\end{center}

\noindent Requires $n \ge 3$ (g57 gates need three distinct wires). Invocation:
\begin{quote}
\code{sss --cnot --sliced-sandwich -n <n> {-}{-}sandwich-m <m> {-}{-}sandwich-s <s>
[{-}{-}sandwich-balanced] \dots}
\end{quote}
\noindent The variant is chosen per run, never per build, and defaults to
\code{classic} everywhere. For \code{gen\_sandwich\_gadget} it is the final
positional argument or the environment variable \code{SANDWICH\_VARIANT}
(an explicit positional wins).
The transformed size before mixing is
$\lvert C\rvert + m + n + 2s$ gates
(source, $D$, copy step, two slice blocks); e.g.\ for $n=8$ with a $40$-gate
source and defaults ($s=24$, $m=72$): $40+72+8+48 = 168$ gates, of which
$40+72=112$ are g57.

\section{Verification}

The driver checks the pinned contract after transformation and after every
mixing round: the \code{SandwichSecond} view for the classic variant (the
second half of the input is forced to zero and the second-half output must
equal $C(x)$) and the \code{SandwichFirst} view for the balanced one (same
pinned slice, first-half output must equal $C(x)$). Unit tests cover
slice-gate deadness on both sides, each variant's forward slice contract with
a permutation check, the classic inverse's second-half bijection on the slice,
the balanced inverse's exact recovery of $D^{-1}$ on the mirrored slice, and
the float stage's band for both variants.

\vspace{6pt}
\noindent\textbf{Provenance.} Code: \code{SandwichVariant},
\code{sliced\_sandwich\_cnot}, \code{sliced\_sandwich\_with\_d},
\code{sliced\_sandwich\_build}, \code{random\_g57\_xgates},
\code{sandwich\_slice\_gates}, \code{shift\_xgate\_wires},
\code{random\_interleave} in \code{src/preprocessing/gadgets.rs}; driver branch and
\code{FunctionView::SandwichSecond} / \code{FunctionView::SandwichFirst} in
\code{src/db\_mixing/main\_mix\_cnot.rs}; flags in \code{src/main.rs} and
\code{src/commands/sss.rs}; variant argument in
\code{src/preprocessing/bin/gen\_sandwich\_gadget.rs}. Classic circuits are
unchanged: the choice bit defaults to \code{classic} and the builder draws the
same values in the same order as before, so a classic seed still reproduces
its circuit bit for bit. Companion quick-reference:
\code{docs/SLICED\_SANDWICH.md}.

\end{document}
